\documentclass[11pt,a4paper]{article}
\usepackage[utf8]{inputenc}
\usepackage[spanish,es-tabla]{babel}
\usepackage{amsmath,amssymb}
\usepackage{booktabs}
\usepackage{geometry}
\geometry{margin=2.5cm}
\usepackage{hyperref}
\usepackage{enumitem}
\usepackage{xcolor}

\hypersetup{
    colorlinks=true,
    linkcolor=blue!70!black,
    urlcolor=blue!70!black
}

\title{\textbf{Plan de Trabajo: Depuración y Preparación de Datos SUI}\\
\large Proyecto E: Clima y Demanda de Energía Eléctrica en el Caribe Colombiano (\texttt{caribbean-climate-demand})\\
\large \textbf{Fase Data Gate E2}}
\author{\textbf{Colaborador:} \texttt{Rafael Marulanda}\\
\textbf{Destinatarios:} Equipo de Investigación (3 Colaboradores)}
\date{\today}

\begin{document}

\maketitle

\begin{abstract}
Este documento establece la hoja de ruta metodológica
para la ingesta, estandarización y depuración de las bases de datos 
del Sistema Unificado de Información (SUI) de la Superintendencia de 
Servicios Públicos Domiciliarios (SSPD). Con el fin de blindar el avance 
científico del proyecto frente a restricciones de acceso público a los datos, 
se adopta una \textbf{estrategia dual en paralelo}: (i) \textbf{Vía A (Régimen CREG 015/2018 / 2021--2024)}, 
enfocada en los nuevos operadores Afinia y Air-e, supeditada a la completitud de lotes nacionales 
solicitados mediante derecho de petición formal ante la SSPD; y (ii) \textbf{Vía B (Régimen CREG 097/2008 / 2018--2020 y serie histórica)}, 
enfocada en el operador integrado \textit{Electricaribe} mediante los formatos históricos 
Formato 2 (SUI 439) y Formato 1 (SUI 438), cuya ejecución es inmediata con datos ya disponibles localmente.
Ambas ramas alimentan los criterios del \textbf{Gate E2}.
\end{abstract}

\tableofcontents
\newpage

\section{Justificación Metodológica: Estrategia Dual de Estimación}

Para asegurar la continuidad empírica y la robustez del análisis econométrico, 
el plan de depuración no descarta ninguna fuente, sino que articula dos líneas de trabajo concurrentes:

\subsection{Vía A: Régimen Post-2020 (CREG 015/2018 / 2021--2024)}
\begin{enumerate}[label=\alph*.]
    \item \textbf{Cobertura y Operadores:} Comprende los mercados de comercialización y distribución 
    de \textbf{Afinia} (CaribeMar de la Costa S.A.S. E.S.P.) y \textbf{Air-e} (Air-e S.A.S. E.S.P.).
    \item \textbf{Estructura Granular:} Se sustenta en los formatos transaccionales modernos:
    \begin{itemize}
        \item \textbf{TC2 (SUI 1743):} Facturación por suscriptor, con fechas exactas de lectura anterior y actual 
        (\texttt{FCH\_LECTURA\_ANT}, \texttt{FCH\_LECTURA\_ACT}) y registro directo de \texttt{TIPO\_LECTURA} (\textit{Real} vs. \textit{Estimada}).
        \item \textbf{TC1 (SUI 1732):} Parámetros técnicos, nivel de tensión, código de transformador y localización geográfica (\texttt{DANE\_NIU}).
    \end{itemize}
    \item \textbf{Estado de Disponibilidad y Gestión Institucional:} Dado que el portal público de datos 
    abiertos únicamente cuenta con los lotes completos del año 2023, la completitud de los años 
    2021, 2022 y 2024 se tramita mediante el \textbf{Derecho de Petición} formal radicado ante la SSPD 
    (documentado en \texttt{docs/derch\_petic.pdf}). La infraestructura de código queda preparada para procesar 
    dichos lotes de forma inmediata tan pronto sean suministrados.
\end{enumerate}

\subsection{Vía B: Régimen Histórico (CREG 097/2008 / 2018--2020 y Serie Extendida)}
\begin{enumerate}[label=\alph*.]
    \item \textbf{Operador Único Integrado:} Cubre los 7 departamentos de la región Caribe bajo la operación de 
    \textbf{Electrificadora del Caribe S.A. E.S.P. (Electricaribe)}, eliminando el quiebre de mercado y la división 
    territorial ocurrida tras la intervención y liquidación en octubre de 2020.
    \item \textbf{Factibilidad Inmediata y Calidad de Medición (Formato 2 / SUI 439):}
    \begin{itemize}
        \item A diferencia de la creencia inicial, la auditoría empírica del encabezado de los microdatos de SUI 439 
        demuestra que \textbf{sí se reporta la variable \texttt{Tipo de Lectura}} (\texttt{R} = Real, \texttt{E} = Estimada, \texttt{P} = Promedio), 
        permitiendo aislar con total rigurosidad las observaciones de medición real.
        \item Asimismo, el Formato 2 contiene directamente el \texttt{Estrato} socioeconómico (1 a 6) y la desagregación 
        geográfica (\texttt{Codigo DANE}, \texttt{Departamento}, \texttt{Municipio}), agilizando la depuración.
        \item \textbf{Reconstrucción de la Exposición Climática:} La ventana de exposición horaria (ERA5-Land) 
        se construye determinísticamente a partir de la fecha de inicio del ciclo ($t_0$) y los días facturados ($D$): 
        $[t_0, \; t_0 + D \text{ días}]$.
    \end{itemize}
    \item \textbf{Extensibilidad Histórica (Pre-2017):} Esta vía permite incorporar años anteriores (ej. 2014--2017) 
    para evaluar eventos climáticos extremos de ciclo largo (como el Fenómeno de El Niño 2015--2016) y proveer 
    un contrafactual limpio no contaminado por el choque de confinamiento COVID-19 de 2020.
\end{enumerate}

---

\section{Estructura de Trabajo para los 3 Colaboradores (Ejecución Paralela)}

La división de responsabilidades se organiza en tres frentes coordinados para atender ambas vías:

\subsection{Colaborador 1: Ingesta, Descargas Masivas y Filtrado Regional Caribe}
\begin{itemize}
    \item \textbf{Línea Vía A (CREG 015 / 2021--2024):}
    \begin{itemize}
        \item Monitorear e incorporar los lotes nacionales solicitados en el derecho de petición indexados en \texttt{data/manifests/sui\_2021\_2024\_manifest.json}.
        \item Filtrar Afinia y Air-e en \texttt{src/data/02\_filter\_caribe\_sui.py} y exportar a \texttt{data/intermediate/sui\_caribe\_tc1\_raw.parquet} y \texttt{sui\_caribe\_tc2\_raw.parquet}.
    \end{itemize}
    \item \textbf{Línea Vía B (CREG 097 / 2018--2020 y Pre-2017):}
    \begin{itemize}
        \item Catalogar los archivos históricos de Formato 2 (SUI 439) y Formato 1 (SUI 438) disponibles en \texttt{data/raw/sui/} (Google Drive \texttt{00\_raw\_sui/}) en el manifiesto \texttt{data/manifests/sui\_2018\_2020\_manifest.json}.
        \item Implementar el filtrado de \textit{Electricaribe} (\texttt{ID\_EMPRESA == '2249'} oficial y \texttt{'2238'}, nombre del operador, o códigos DANE del Caribe: Atlántico, Bolívar, Cesar, Córdoba, La Guajira, Magdalena, Sucre).
        \item Exportar los extractos regionales históricos a \texttt{data/intermediate/sui\_caribe\_f1\_raw.parquet} (77,4 millones de registros) y particiones anuales \texttt{sui\_caribe\_f2\_2018.parquet}, \texttt{sui\_caribe\_f2\_2019.parquet} y \texttt{sui\_caribe\_f2\_2020.parquet} (54,8 millones de transacciones).
    \end{itemize}
\end{itemize}

\subsection{Colaborador 2: Depuración de Facturación, Ciclos y Consumo}
\begin{itemize}
    \item \textbf{Consistencia de Fechas y Ciclos Mensuales:} 
    Restringir la muestra primaria a ciclos estándar de facturación mensual en ambas vías:
    $$\text{\texttt{Dias Facturados}} \in [25, 35]$$
    \item \textbf{Filtrado Estricto de Calidad de Lectura:}
    \begin{itemize}
        \item En Vía A (TC2): Exigir estrictamente $\text{\texttt{TIPO\_LECTURA}} == \text{'REAL'}$.
        \item En Vía B (Formato 2): Exigir estrictamente $\text{\texttt{Tipo de Lectura}} == \text{'R'}$.
        \item Aislar y reportar atrición de lecturas estimadas o promediadas en \texttt{src/data/03\_clean\_sui.py}.
    \end{itemize}
    \item \textbf{Construcción de la Variable Dependiente:} Calcular el consumo diario promedio normalizado para cada observación suscriptor-mes ($Y_{it}$ en kWh/día):
    $$Y_{it} = \frac{\text{\texttt{Consumo}}_{it}}{\text{\texttt{Dias Facturados}}_{it}}$$
    \item \textbf{Homologación del Intervalo Climático:} Generar para cada suscriptor el par ordenado de fechas $[t_{\text{inicio}}, t_{\text{fin}}]$ para acoplar la acumulación de grados-día de refrigeración (CDD) de ERA5-Land.
    \item \textbf{Salida Intermediate:} Generar \texttt{data/intermediate/sui\_caribe\_tc2\_clean.parquet} (Vía A) y \texttt{data/intermediate/sui\_caribe\_f2\_clean.parquet} (Vía B).
\end{itemize}

\subsection{Colaborador 3: Limpieza Catastral, Estratificación y Red Eléctrica}
\begin{itemize}
    \item \textbf{Estabilidad Longitudinal de la Llave NIU:} Verificar la permanencia e invariancia temporal del 
    Número de Identificación del Usuario (\texttt{NIU}) dentro del operador correspondiente.
    \item \textbf{Sectorización Residencial:} Filtrar suscriptores de uso \textbf{Residencial} (\texttt{Estrato} 1 a 6), descartando sector comercial, industrial o provisional.
    \item \textbf{Homologación Geográfica DANE:} Estandarizar la codificación municipal (5 dígitos) y de centros poblados (8 dígitos) para los 7 departamentos del Caribe.
    \item \textbf{Topología de Red y GIS:} Cruzar identificadores de circuitos y transformadores de distribución con los datos espaciales UPME en \texttt{data/spatial/}.
    \item \textbf{Salida Intermediate:} Generar \texttt{data/intermediate/sui\_caribe\_tc1\_clean.parquet} (Vía A) y \texttt{data/intermediate/sui\_caribe\_f1\_clean.parquet} (Vía B).
\end{itemize}

---

\section{Cronograma de Entregables (Ejecución Paralela)}

\begin{table}[h!]
\centering
\caption{Cronograma de Hitos de Depuración y Ensamble SUI (Régimen Dual)}
\label{tab:cronograma}
\begin{tabular}{@{}lllp{7.2cm}@{}}
\toprule
\textbf{Hito} & \textbf{Plazo} & \textbf{Responsable} & \textbf{Entregable Concreto} \\ \midrule
H1 & Semana 1 & Colaborador 1 & Extracción regional de Formato 2 y Formato 1 (2018--2020 / histórico) a Parquet; manifiesto con hashes; monitoreo de respuesta a derecho de petición 2021--2024. \\
H2 & Semana 2 & Colaborador 2 & Base depurada de facturación (\texttt{sui\_caribe\_f2\_clean.parquet} y \texttt{tc2\_clean}) con lecturas reales y consumo diario normalizado. \\
H3 & Semana 2 & Colaborador 3 & Base de atributos de usuarios (\texttt{sui\_caribe\_f1\_clean.parquet} y \texttt{tc1\_clean}) con estrato validado y códigos DANE. \\
H4 & Semana 3 & Equipo Conjunto & Agregación climática ERA5-Land adaptada al intervalo de ciclo e integración de tarifas CREG CU. \\
H5 & Semana 3 & \texttt{econ-research-director} & Matriz final analítica de estimación (\texttt{panel\_caribe\_2018\_2020\_v1.parquet} y esquema complementario 2021--2024) para revisión y aprobación \textbf{Gate E2}. \\ \bottomrule
\end{tabular}
\end{table}

---

\section{Criterios de Aceptación para el Gate E2}

Para autorizar la liberación formal de datos (\textbf{Gate E2}), se debe verificar:
\begin{enumerate}
    \item \textbf{Trazabilidad Total:} Todo registro del panel final debe rastrearse a un archivo zip original registrado en los manifiestos JSON con checksum SHA-256.
    \item \textbf{Calidad de Medición:} El 100\% de las observaciones del panel primario deben provenir de lecturas reales (\texttt{TIPO\_LECTURA == 'REAL'}).
    \item \textbf{Reporte de Atrición Transparente:} Matriz explícita de pérdidas de observaciones por filtros aplicados (empresa, días facturados, lecturas estimadas, valores atípicos).
    \item \textbf{Independencia Climática:} La muestra debe depurarse por completo \textbf{antes} de cruzar con datos de temperatura (ERA5-Land) para prevenir sesgo de selección o \textit{p-hacking}.
\end{enumerate}

\end{document}
